% !TEX root = Séquence 2 - Analyse didactique - Activité 3 Somme de vecteurs.tex
% Document de travail personnel (non destiné aux élèves) : analyse
% didactique et pédagogique de l'Activité 3 "Enchaîner deux translations"
% (voir "Séquence 2 - Activités - Vecteurs.tex"), écrite avant la mise en
% oeuvre en classe. Reprend le montage "article" simple déjà utilisé pour
% "Séquence 1 - Travail d'analyse des exercices.tex" plutôt que
% CLT2627.sty, qui est pensé pour les documents élèves.
\documentclass[11pt,a4paper]{article}
\usepackage[utf8]{inputenc}
\usepackage[T1]{fontenc}
\usepackage[french]{babel}
\renewcommand{\familydefault}{\sfdefault}
\usepackage[left=2cm,right=2cm,top=2cm,bottom=2cm]{geometry}
\usepackage{amsmath,amssymb}
\usepackage{xcolor}
\usepackage{titlesec}
\usepackage{enumitem}
\usepackage{array}
\usepackage{tcolorbox}
\tcbuselibrary{skins,breakable}
\usepackage[hidelinks]{hyperref}

\definecolor{infoblue}{HTML}{1C7ED6}
\definecolor{warnorange}{HTML}{E8590C}

\titleformat{\section}{\Large\bfseries\color{black}}{\thesection.}{6pt}{}[{\color{gray!40}\titlerule}]
\titleformat{\subsection}{\large\bfseries}{\thesubsection}{6pt}{}
\titlespacing*{\section}{0pt}{16pt}{8pt}
\setlength{\parindent}{0pt}

\newtcolorbox{Vigilance}[1][]{enhanced,breakable,colback=warnorange!6,colframe=warnorange,
  boxrule=0.8pt,arc=2pt,left=8pt,right=8pt,top=5pt,bottom=5pt,
  title={\bfseries Point de vigilance},#1}
\newtcolorbox{Lien}[1][]{enhanced,breakable,colback=infoblue!6,colframe=infoblue,
  boxrule=0.8pt,arc=2pt,left=8pt,right=8pt,top=5pt,bottom=5pt,
  title={\bfseries Lien avec la suite du cours},#1}

% \Vj n'existe pas dans ce document autonome (pas de CLT2627.sty) : on
% redéfinit ici la notation vecteur utilisée dans "Séquence 2 - Activités
% - Vecteurs.tex", pour citer les mêmes objets ($\Vj u$, $\Vj{AB}$...)
% avec la même notation que le document source. Défini dans le préambule
% (et non après \maketitle) : utilisé dès le tableau d'en-tête.
\newcommand{\Vj}[1]{\overrightarrow{#1}}

\title{\textbf{Analyse didactique et pédagogique}\\[4pt]
\large Activité  --- \og Enchaîner deux translations \fg{}}
\author{Séquence 2 --- Vecteurs du plan}
\date{}

\begin{document}
\maketitle
\vspace{-1.5em}

\noindent
\begin{tabular}{@{}ll}
\textbf{Niveau :} & Secondes \\
\textbf{Durée indicative :} & 15 min \\
\textbf{Place dans la progression :} & après les pages 1--2 du cours (vecteur $\Vj{AB}$, translation associée,\\
 & égalité de vecteurs, vecteur nul/opposé) ; avant la section \og Somme de\\
 & deux vecteurs \fg{} (définition, méthode bout à bout, relation de Chasles).
\end{tabular}

\section{Objectif d'apprentissage visé}

Faire émerger, par la manipulation plutôt que par un exposé, l'idée que
\textbf{enchaîner deux translations de vecteurs $\Vj{u}$ et $\Vj{w}$
équivaut à une seule translation de vecteur $\Vj{u}+\Vj{w}$}, et préparer
intuitivement la méthode \og bout à bout \fg{} que le cours formalisera
juste après. L'activité ne vise pas à \emph{définir} la somme (ce sera
fait dans le cours), mais à donner du sens à cette définition avant
qu'elle n'arrive.

\section{Analyse de la situation-problème}

\begin{itemize}[leftmargin=1.4em,itemsep=6pt, label=$\star$]
\item \textbf{Contexte réaliste et dynamique (hockey)} : le vecteur est
vécu comme un déplacement du palet, dans un registre concret, avant son
institutionnalisation symbolique en cours. Bon levier de motivation pour
une première approche.

\item \textbf{Réinvestissement du vocabulaire de représentant} : au
moment de placer $B$, c'est l'occasion de réinvestir au
tableau le vocabulaire installé en cours sur les représentants d'un
vecteur --- formuler par exemple à l'oral \og on trace un représentant
d'origine $A$ de $\Vj{u}$ \fg{}, plutôt que de rester au seul registre
familier de la \og passe \fg{}. Ce geste relie explicitement la
construction (jeu de hockey) à la notion formelle déjà installée, et
prépare le passage au registre symbolique qui suivra en question d).

\item \textbf{Vecteur résultat déjà présent sur la figure} (le vecteur
$\Vj{n}$ est tracé dès le départ, parmi $\Vj{m}$, $\Vj{v}$, $\Vj{p}$) :
c'est la variable didactique clé de l'activité. Elle transforme la tâche
de \og construire bout à bout \fg{} en tâche de \textbf{reconnaissance
perceptive} (comparer direction/sens/norme avec les vecteurs déjà
tracés).
\begin{itemize}[leftmargin=1.2em,itemsep=3pt, label=$\odot$]
\item \emph{Avantage} : rend la conjecture accessible à toute la classe,
sans risque qu'une erreur de tracé bloque la découverte.
\item \emph{Point de vigilance} : un élève peut repérer $\Vj{n}$ par
élimination visuelle sans avoir vraiment compris le mécanisme de
chaînage --- prévoir une question orale de contrôle du type \og comment
as-tu vérifié que c'est bien ce vecteur-là et pas un autre ? \fg{}
(retour aux critères d'égalité de vecteurs vus en page 2 : même
direction, même sens, même norme).
\end{itemize}

\item \textbf{Réinvestissement immédiat d'un acquis récent} : pour
identifier $\Vj{n}$, l'élève mobilise exactement le critère d'égalité de
deux vecteurs installé en page 2. Bonne articulation spiralaire --- un
savoir tout juste stabilisé devient l'outil qui permet de construire le
suivant.

\item \textbf{Institutionnalisation dosée} : l'encadré jaune de
l'activité donne le nom \og somme \fg{} juste après la construction,
mais \emph{avant} que l'élève n'ait fini d'identifier $\Vj{n}$. Ce séquencement sépare bien deux moments cognitifs :
(1) construire/observer, (2) nommer --- cohérent avec une démarche
d'investigation guidée plutôt qu'une découverte totalement autonome. En
vidéoprojection, il peut être utile de masquer cet encadré jusqu'à ce
que la classe ait elle-même formulé une conjecture, pour ne pas la
\og griller \fg{}.

\newpage 

\item \textbf{Deuxième question} ($\Vj{p}+\Vj{v}$, points D et E) :
rejouée sans filet, elle sert de
transfert en autonomie immédiate après l'institutionnalisation --- bonne
gradation de l'étayage.

\end{itemize}

\begin{Vigilance}
L'énoncé insiste : \og un seul déplacement final \fg{} --- c'est le
message le plus important de l'activité, car il prépare directement
l'idée d'\textbf{indépendance du chemin suivi}, qui est le coeur de la
relation de Chasles vue juste après dans le cours. Il vaut la peine de
l'expliciter à l'oral en mise en commun plutôt que de le laisser filer.
\end{Vigilance}

\section{Scénario de séance possible (15 min)}

\begin{center}
\renewcommand{\arraystretch}{1.3}
\begin{tabular}{|>{\raggedright\arraybackslash}p{2.2cm}|p{12.5cm}|}
\hline
\textbf{Temps approximatif} & \textbf{Déroulé} \\
\hline
2 min & Lecture de la consigne, contextualisation (rappel : translation = vecteur associé, déjà vu) \\
\hline
5 min & Recherche collective : placer B, C, tracer le vecteur unique --- pas de calcul, uniquement du tracé \\
\hline
3 min & Mise en commun : un élève identifie $\Vj{n}$, le professeur révèle/valide l'encadré de synthèse \\
\hline
4 min & 2\ieme{} sous-question (D, E, $\Vj{p}+\Vj{v}$) en autonomie, transfert immédiat \\
\hline
 & Clôture $\to$ enchaînement sur le cours \S 2 (définition + méthode bout à bout)\\
\hline
\end{tabular}
\end{center}

\begin{Lien}
Le protocole vécu ici (A $\to$ B via $\Vj{u}$, puis B $\to$ C via
$\Vj{w}$) est très exactement celui que formalisera le cours page 3
(\og on choisit un point M, on construit N tel que $\Vj{MN}=\Vj{u}$,
puis P tel que $\Vj{NP}=\Vj{v}$ \fg{}), à ceci près que dans l'activité
le point de départ est imposé (A) alors que le cours généralise à un
point M quelconque --- simplification pertinente pour une première
rencontre. L'activité fonctionne donc bien comme un \textbf{dispositif
de découverte qui précède et motive l'institutionnalisation}, plutôt
qu'une application du cours.
\end{Lien}

\end{document}
